\chapter{Appendix D: Developer Notes (Extended)}

This appendix expands on the short Developer Notes section at the front of the book.
It\textquotesingle s written as an engineering log: what I built, why, and what I learned.

\section{From \textit{reading} to \textit{building}}
A surprising part of framework engineering is that you can understand the concepts intellectually and still struggle to implement them.
A renderer forces precision.
A scheduler forces you to choose semantics.
A compiler forces you to state what is safe.

The starting point for my learning journey was the book:

\begin{quote}
	extit{Build a Frontend Web Framework (From First Principles)}\\
\textsc{ÁNGEL SOLA ORBAICETA}
\end{quote}

It is a solid introduction to the core loop of building a small framework.
I used it to bootstrap the baseline mental model, then I iterated.

\subsection{How to use the inspiration book alongside this repository}
An approach that works well:

\begin{enumerate}
  \item Implement the baseline renderer and component model from \textit{Build a Frontend Web Framework (From First Principles)}.
  \item Then open this repository and locate the matching unit tests under \texttt{src/\_\_tests\_\_/}.
  \item Port those tests to your own project.
  \item Only then start optimizing (diffing, scheduling, HRBR), keeping the tests unchanged.
\end{enumerate}

This forces a clean separation between \emph{behavior} (what must stay true) and \emph{implementation} (what you are allowed to change).

\section{What \texttt{\_\_tests\_\_} taught me}
The test suite is the highest-leverage part of this project.
Tests do three jobs simultaneously:

\begin{itemize}
  \item They define the contract (what behavior counts as \emph{correct}).
  \item They prevent accidental regressions.
  \item They document edge cases that are easy to forget.
\end{itemize}

For example, lifecycle timing is not obvious when you have async scheduling.
It only becomes obvious when a test waits for \texttt{nextTick()} and asserts an effect has run exactly once.

\paragraph{Verification discipline (tests are not an afterthought)}
When this book says \textquotedblleft tests are executable specifications\textquotedblright, it means:

\begin{itemize}
  \item every important claim should be traceable to an expectation
  \item every bug fix should ship with a regression test
\end{itemize}

It\textquotesingle s slower on day one, but dramatically faster by the time your framework has a scheduler, a diff algorithm, and a compiler.

\section{A short author bio (and why performance shows up so often)}
I\textquotesingle m primarily a backend and systems engineer.
I\textquotesingle ve worked on performance-sensitive systems (including quantitative infrastructure), so I\textquotesingle m naturally drawn to:

\begin{itemize}
  \item deterministic semantics (especially around scheduling),
  \item careful measurement (benchmarks before and after),
  \item and packaging contracts (exports maps that match reality).
\end{itemize}

That bias is part of why Relax.js tries to be a \emph{real library} and not only a demo.

\section{Why algorithms appear in UI code}
A framework looks like \textquotedblleft DOM glue\textquotedblright until you hit lists.
Then you need algorithms:

\begin{itemize}
  \item matching old children to new children
  \item detecting moves vs inserts vs deletes
  \item minimizing expensive DOM operations
\end{itemize}

That\textquotesingle s why this repository contains diff/sequence logic and tests to keep it honest.
Algorithmic thinking is not optional if you care about correctness and scalability.

\section{Why compiler assistance is worth it}
The HRBR path exists because many UI trees are structurally stable.
If structure is stable, you can treat DOM updates as \textquotedblleft update slots\textquotedblright rather than \textquotedblleft diff a tree\textquotedblright.

But the key word is \emph{safe}.
The compiler must be conservative.
When it cannot prove stability (conditionals that change structure, dynamic tag names, etc.) it emits a fallback.
That fallback has to be a first-class, tested path.

\section{Engineering principle: make your fast path debuggable}
Fast paths become liabilities when they are invisible.
That\textquotesingle s why devtools/instrumentation exists.
If a performance win cannot be measured and explained, it will eventually be removed (or, worse, it will hide bugs).

\section{Where to go next}
If you want to evolve Relax.js further (or build your own):

\begin{itemize}
  \item Add a small router and prove it doesn\textquotesingle t leak listeners.
  \item Add async data fetching primitives with cancellation.
  \item Explore hydration and partial hydration trade-offs.
  \item Add profiling hooks around patch phases (per component / per block).
\end{itemize}
